% ==============================================================================
%  2.2  El modelo de regresión lineal simple: formulación, hipótesis e
%       interpretación de los parámetros
% ==============================================================================
\section{El modelo de regresión lineal simple}
\label{sec:t2-modelo}

El modelo de regresión más sencillo es el \textbf{modelo de regresión lineal simple}, en el
que intervienen:
\begin{itemize}
  \item una variable respuesta $Y$ cuantitativa;
  \item una única variable explicativa $X$ cuantitativa.
\end{itemize}
Se estudia si los cambios en $X$ afectan \emph{de forma lineal} a $Y$ ($Y\leftarrow X$).

\begin{ejemplo}[Conductores: formulación del modelo][ej:t2-conductores]
  En el \cref{ej:conductores} (Tema 1) la edad es la variable explicativa $X$ y la distancia
  máxima de lectura de la señal es la respuesta $Y$. El diagrama de dispersión sugiere una
  relación lineal negativa, por lo que se modeliza la distancia media como una función lineal
  de la edad más un error aleatorio, que es el modelo que se define a continuación.
\end{ejemplo}

\subsection{Formulación del modelo}

\begin{definicion}[Modelo de regresión lineal simple][def:rls]
  \[
    Y_i = \beta_0 + \beta_1 X_i + \varepsilon_i, \qquad i=1,\ldots,n,
  \]
  donde
  \begin{itemize}
    \item $Y_i$ es la respuesta del $i$-ésimo individuo;
    \item $\beta_0$ y $\beta_1$ son \textbf{parámetros}, llamados \textbf{coeficientes de la
      regresión};
    \item $X_i$ es el valor de la variable explicativa (predictor) del $i$-ésimo individuo.
      \textbf{Se asume conocido} (no aleatorio);
    \item $\varepsilon_i$ representa la componente aleatoria: es una \textbf{variable
      aleatoria} $\Normal(0,\sigma^2)$, y $\varepsilon_1,\ldots,\varepsilon_n$ son
      \textbf{i.i.d.}
  \end{itemize}
\end{definicion}

\paragraph{Parámetros del modelo.}
\begin{itemize}
  \item $\beta_0$: \textbf{\emph{intercept}}, u ordenada en el origen.
  \item $\beta_1$: \textbf{pendiente}. $\beta_1=0$ si no hay relación (lineal) entre $X$ e $Y$.
  \item $\sigma^2$: \textbf{varianza del error}.
\end{itemize}

\subsection{Hipótesis del modelo}
\label{sec:t2-hipotesis}

Al ajustar un modelo de regresión lineal simple se asume que los \textbf{errores} son variables
aleatorias independientes e igualmente distribuidas, \textbf{normales con media 0 y varianza
constante}:
\[
  \varepsilon_i \sim \Normal(0,\sigma^2) \ \text{independientes}, \qquad \forall i=1,\ldots,n.
\]
Esto supone cuatro hipótesis, que deben comprobarse para que el modelo sea válido:
\begin{enumerate}
  \item \textbf{Linealidad}: $\E(\varepsilon_i)=0$ (la media de $Y$ es realmente la recta).
  \item \textbf{Homogeneidad de varianzas} (homocedasticidad): $\Var(\varepsilon_i)=\sigma^2$,
    que no depende de $i$.
  \item \textbf{Normalidad}: $\varepsilon_i \sim \Normal$.
  \item \textbf{Independencia}: $\varepsilon_i$ y $\varepsilon_j$ independientes,
    $\forall i\neq j$.
\end{enumerate}

\begin{proposicion}[Distribución de la respuesta][prop:t2-dist-Y]
  Las hipótesis sobre $\varepsilon$ son equivalentes a
  \[
    Y_i \sim \Normal(\beta_0+\beta_1X_i,\ \sigma^2) \ \text{independientes},
    \qquad \forall i=1,\ldots,n .
  \]
\end{proposicion}

En efecto, la respuesta $Y_i=\beta_0+\beta_1X_i+\varepsilon_i$ es la suma de una constante,
$\beta_0+\beta_1X_i$ (los $X_i$ son conocidos), y de la variable aleatoria $\varepsilon_i$:
\begin{itemize}
  \item $Y_i$ es una variable aleatoria normal, por ser combinación lineal de variables
    aleatorias normales;
  \item $\E(Y_i)=\beta_0+\beta_1X_i$, ya que $\E(\varepsilon_i)=0$;
  \item $\Var(Y_i)=\Var(\varepsilon_i)=\sigma^2$, ya que sumar una constante no altera la
    varianza;
  \item las respuestas de individuos distintos son incorreladas: para $i\neq j$,
    \[
      \Cov(Y_i,Y_j)=\Cov(\beta_0+\beta_1X_i+\varepsilon_i,\ \beta_0+\beta_1X_j+\varepsilon_j)
      =\Cov(\varepsilon_i,\varepsilon_j)=0,
    \]
    por la independencia de los errores. Además, $Y_i$ e $Y_j$ son independientes, ya que cada
    una es función únicamente de su propio error y $\varepsilon_i$ y $\varepsilon_j$ son
    independientes.
\end{itemize}

\begin{nota}
  Como los errores $\varepsilon_i$ no se observan, las hipótesis se comprueban sobre los
  \emph{residuos} $e_i$ (\cref{def:t2-ajuste}), mediante el \textbf{análisis residual}.
\end{nota}

\subsection{Interpretación de los parámetros}

La parte determinista del modelo, $\beta_0+\beta_1x_i$, especifica que la \textbf{media
poblacional} de la respuesta $Y$ cuando $X=x_i$ es una \textbf{línea recta}:
\[
  Y_i = \underbrace{\beta_0+\beta_1x_i}_{\E(Y_i\cond X=x_i)\,=\,\E(Y_i)} + \varepsilon_i .
\]
Por ser una recta:
\begin{itemize}
  \item $\beta_0$ (\textbf{\emph{intercept}}) es el valor esperado de la respuesta cuando $X=0$;
  \item $\beta_1$ (\textbf{pendiente}) es el incremento (o decremento, si es negativa) de la
    media de la respuesta cuando $X$ aumenta 1 unidad.
\end{itemize}

\begin{figure}[H]
  \centering
  \begin{tikzpicture}[x=0.11cm, y=0.13cm]
    \draw[-{Stealth}] (0,0) -- (70,0) node[below] {$X$};
    \draw[-{Stealth}] (0,0) -- (0,40) node[left] {$Y$};
    \foreach \t in {10,20,30,40,50,60} \draw (\t,0) -- (\t,-0.6) node[below, font=\scriptsize] {\t};
    \foreach \t in {10,20,30} \draw (0,\t) -- (-0.8,\t) node[left, font=\scriptsize] {\t};
    \draw[NavyBlue, thick] (0,0) -- (68,34) node[right, font=\scriptsize] {$\E(Y)=0+0.5X$};
    \foreach \xx/\lab/\dy in {20/1/2.2,30/3/7,55/2/2.2} {
      \pgfmathsetmacro{\mm}{0.5*\xx}
      \draw[gray] (\xx,{\mm-8}) -- (\xx,{\mm+8});
      \draw[colRes, thick, domain=-8:8, samples=40] plot ({\xx+4.5*exp(-(\x)^2/10)}, {\mm+\x});
      \fill[NavyBlue] (\xx,\mm) circle (1.3pt);
      \node[font=\scriptsize, anchor=east] at (\xx-0.8,\mm+\dy) {$\E(Y_{\lab})=\pgfmathprintnumber{\mm}$};
    }
    \fill (20,7.3) circle (1pt) node[below right, font=\scriptsize] {$Y_1$};
    \fill (55,29.5) circle (1pt) node[below right, font=\scriptsize] {$Y_2$};
    \fill (30,16.4) circle (1pt) node[right, font=\scriptsize] {$Y_3$};
  \end{tikzpicture}
  \caption{Recta de regresión poblacional $\E(Y)=0+0.5X$ y distribución normal de $Y$ para tres
    valores de $X$, todas con la misma varianza. Cada observación $Y_i$ es un valor de su
    distribución.}
\end{figure}

\subsection{Ajustar el modelo: recta estimada, valores ajustados y residuos}

Ajustar el modelo consiste en estimar sus parámetros a partir de los datos: la \textbf{recta de
regresión} ($\beta_0$ y $\beta_1$, que describen el modelo) y $\sigma^2$, que permite cuantificar
el error que se comete.

\begin{definicion}[Recta estimada, valores ajustados y residuos][def:t2-ajuste]
  Sean $\est{\beta}_0$, $\est{\beta}_1$ y $\est{\sigma}^2$ los estimadores de los parámetros.
  \begin{itemize}
    \item La \textbf{recta de regresión estimada} es $\est{Y}=\est{\beta}_0+\est{\beta}_1X$.
    \item Los \textbf{valores ajustados} son $\est{y}_i=\est{\beta}_0+\est{\beta}_1x_i$,
      $i=1,\ldots,n$: los valores estimados por el modelo.
    \item Las desviaciones respecto de la verdadera recta son los \textbf{errores}
      $\varepsilon_i=y_i-(\beta_0+\beta_1x_i)$, que son desconocidos porque también lo son los
      coeficientes. Se estiman mediante los \textbf{residuos}:
      \[
        e_i = y_i-\est{y}_i = y_i-(\est{\beta}_0+\est{\beta}_1x_i).
      \]
  \end{itemize}
\end{definicion}

\begin{figure}[H]
  \centering
  \begin{tikzpicture}
    \begin{axis}[width=8cm, height=5.5cm, axis lines=left, xtick={4.5}, xticklabels={$x_i$},
        ytick={3.2,4.55}, yticklabels={$y_i$,$\est{y}_i$}, xmin=0, xmax=10, ymin=0, ymax=9.5,
        clip=false]
      \addplot[recta, domain=0.3:9.8] {1.4+0.7*x};
      \node[font=\scriptsize, NavyBlue, anchor=west] at (axis cs:9.9,8.3)
        {$\est{y}=\est{\beta}_0+\est{\beta}_1x$};
      \addplot[puntos, mark size=1.5pt] coordinates {(1.5,3.5) (2,2.7) (2.5,2.8) (3.3,3.4)
        (3.7,3.3) (5.3,5.4) (6.2,6.9) (6.8,5.8) (7.8,6.7) (8.6,8.0) (9.3,7.7)};
      \addplot[only marks, mark=*, mark size=1.8pt, red] coordinates {(4.5,3.2)};
      \addplot[only marks, mark=square, mark size=1.8pt, NavyBlue] coordinates {(4.5,4.55)};
      \draw[dashed, gray] (axis cs:0,3.2) -- (axis cs:4.5,3.2);
      \draw[dashed, gray] (axis cs:0,4.55) -- (axis cs:4.5,4.55);
      \draw[dashed, gray] (axis cs:4.5,0) -- (axis cs:4.5,3.2);
      \draw[red, thick] (axis cs:4.5,3.2) -- (axis cs:4.5,4.55)
        node[midway, right, font=\scriptsize] {$e_i$ (residuo)};
      \node[font=\scriptsize, red, anchor=north west] at (axis cs:4.6,3.1)
        {$(x_i,y_i)$ observación};
      \node[font=\scriptsize, NavyBlue, anchor=south east] at (axis cs:4.4,4.6)
        {$(x_i,\est{y}_i)$ valor ajustado};
    \end{axis}
  \end{tikzpicture}
  \caption{Observación, valor ajustado y residuo.}
\end{figure}

Los estimadores tienen una distribución concreta, lo que permite utilizar \textbf{métodos de
inferencia} para evaluar la significación y la precisión de las estimaciones (Tema 3).

\begin{ejemplo}[Conductores: recta ajustada e interpretación][ej:t2-conductores-ajuste]
  En el ejemplo de los conductores, ajustar el modelo es estimar la recta que representa el
  patrón general de los datos. Se obtiene
  \[
    \widehat{\text{distancia}} = 576 - 3\cdot\text{edad},
    \qquad \est{\beta}_0=576,\quad \est{\beta}_1=-3 .
  \]
  \begin{center}
    \begin{tikzpicture}
      \begin{axis}[width=7.5cm, height=5cm, xlabel={Edad (años)}, ylabel={Distancia (pies)},
          xmin=15, xmax=85, ymin=260, ymax=610, label style={font=\small},
          tick label style={font=\footnotesize}]
        \addplot[puntos, mark size=1.3pt] table[x=edad, y=distancia] {datos/conductores.dat};
        \addplot[red, thick, no marks, domain=17:83] {576-3*x};
        \draw[NavyBlue, -{Stealth}] (axis cs:60,260) -- (axis cs:60,396) -- (axis cs:15,396);
        \node[font=\scriptsize, red, anchor=south west] at (axis cs:40,540)
          {$\widehat{\text{dist.}}=576-3\,\text{edad}$};
      \end{axis}
    \end{tikzpicture}
  \end{center}
  \begin{itemize}
    \item \textbf{Predicción}: para un conductor de 60 años, $576-3\cdot60=396\approx400$ pies.
    \item \textbf{\emph{Intercept}} $\est{\beta}_0=576$: la función de regresión vale 576 en $X=0$,
      lo que equivaldría a decir que un conductor de 0 años vería la señal a 576 pies. En este
      caso \textbf{no es interpretable}, ya que la edad no puede ser 0 (y además queda muy
      fuera del rango observado, 18--82 años).
    \item \textbf{Pendiente} $\est{\beta}_1=-3$: por cada año más de edad, la distancia máxima a
      la que se lee la señal \emph{disminuye, en media}, 3 pies.
  \end{itemize}
\end{ejemplo}

\subsection{Versión centrada del modelo}

El modelo también puede formularse utilizando como predictor $X_i-\media{X}$ (desviación del
valor de $X$ en la observación $i$ respecto de su media) en lugar de $X_i$:
\[
  Y_i = \beta_0^* + \beta_1(X_i-\media{X}) + \varepsilon_i,
  \qquad \text{con } \beta_0^* = \beta_0+\beta_1\media{X}.
\]

\begin{demostracion}[Comprobación]
  Sumando y restando $\beta_1\media{X}$ en el modelo original,
  \begin{align*}
    Y_i &= \beta_0 + \beta_1X_i + \varepsilon_i \\
        &= \beta_0 + \beta_1\media{X} - \beta_1\media{X} + \beta_1X_i + \varepsilon_i \\
        &= \underbrace{(\beta_0 + \beta_1\media{X})}_{\beta_0^*}
          + \beta_1(X_i-\media{X}) + \varepsilon_i .
  \end{align*}
  Ambos modelos representan la misma recta, expresada respecto de un origen distinto en el
  eje $X$.
\end{demostracion}

La interpretación de los parámetros es la siguiente:
\begin{itemize}
  \item $\beta_0^*$ (\emph{intercept}) es el valor esperado de la respuesta cuando $X=\media{X}$.
  \item $\beta_1$ (pendiente) es el cambio en la media de la respuesta cuando $X$ se desvía de
    su media 1 unidad. Es \textbf{el mismo parámetro} que en el modelo con predictor $X_i$.
\end{itemize}

\begin{nota}
  En el ejemplo de los conductores el \emph{intercept} $\beta_0$ no era interpretable, mientras que
  $\beta_0^*$ sí lo es: representa la distancia media para un conductor de edad media.
\end{nota}

\subsection{El proceso de modelización}

\begin{figure}[H]
  \centering
  \begin{tikzpicture}[node distance=0.55cm,
      paso/.style={draw, fill=gray!15, minimum width=4cm, minimum height=0.8cm, align=center,
                   font=\small},
      ext/.style={draw, rounded corners=8pt, fill=gray!30, minimum width=2cm,
                  minimum height=0.7cm, font=\small},
      dec/.style={draw, diamond, aspect=2.2, fill=gray!15, align=center, font=\small,
                  inner sep=1pt},
      fl/.style={-{Stealth}}]
    \node[ext] (ini) {Comienzo};
    \node[paso, below=of ini] (eda) {Análisis exploratorio de datos};
    \node[paso, below=of eda] (aju) {Ajuste de posibles\\ modelos de regresión};
    \node[dec, below=of aju] (ok) {¿Se ajustan bien\\ a los datos?};
    \node[paso, left=1.2cm of ok, minimum width=3.3cm] (rev) {Revisar los modelos\\ y plantear nuevos};
    \node[paso, below=of ok] (mej) {Elegir el ``mejor''};
    \node[paso, below=of mej] (inf) {Hacer inferencias\\ con esos modelos};
    \node[ext, below=of inf] (fin) {Fin};
    \draw[fl] (ini) -- (eda); \draw[fl] (eda) -- (aju); \draw[fl] (aju) -- (ok);
    \draw[fl] (ok) -- node[above, font=\small] {No} (rev);
    \draw[fl] (rev) |- ($(aju.south)!0.5!(ok.north)$);
    \draw[fl] (ok) -- node[right, font=\small] {Sí} (mej);
    \draw[fl] (mej) -- (inf); \draw[fl] (inf) -- (fin);
  \end{tikzpicture}
  \caption{Etapas de un análisis de regresión (compárese con las etapas del ajuste de un modelo
    estadístico, \cref{sec:t1-modelos}).}
\end{figure}
